\section{Detecting periods}

Based on the support of the Fourier transform, one can derive the periods of a function.
Periods are indices $(\thirdcatvariableof{0},\ldots,\thirdcatvariableof{\catorder-1})$ such that
\begin{align*}
    \hypercoreat{\shortcatvariables=\shortcatindices}
    = \hypercoreat{\shortcatvariables=\shortcatindices\oplus\thirdcatindexof{[\catorder]}}
\end{align*}
where $\oplus$ is the tuple position wise sum $\mathrm{mod}\,\,\catdimof{\catenumerator}$.


Each supported index of the Fourier transform corresponds with periods in the function.
The periods of the functions are thus at least the intersection of those (when assuming more structure one might have more periods).

\subsection{Walsh-Hadamard transform}

The Walsh-Hadamard transform is the Fourier transform with respect to the group $\bigtimes_{\catenumeratorin}[2]$.
It is the legwise application of the Hadamard matrix.

In this case, we can easily determine the periods corresponding to each basis vector.
The inverse Fourier transform of $\onehotmapofat{\selindexof{[\catorder]}}{\selvariableof{[\catorder]}}$ is the tensor
\begin{align*}
    \left(\bigotimes_{\catenumeratorin \wcols \selindexof{\catenumerator}=0} \hplusbasat{\catvariableof{\catenumerator}} \right)
    \otimes
    \left(\bigotimes_{\catenumeratorin \wcols \selindexof{\catenumerator}=1} \hminusbasat{\catvariableof{\catenumerator}} \right) \, .
\end{align*}

For each $\hplusbasat{\catvariableof{\catenumerator}}$ we have a period generator of $\onehotmapof{\catenumerator}$.
For each pair of $\hminusbasat{\catvariableof{\catenumerator}}\otimes \hminusbasat{\catvariableof{\seccatenumerator}}$ we have a period generator of $\onehotmapof{\catenumerator}+\onehotmapof{\seccatenumerator}$.

One example is the Deutsch-Jozsa test, where the transformed is the $\onehotmapofat{0}{\selvariableof{[\catorder]}}$ state, if and only if the function is constant.
In this case, any vector is a period, since the group itself is generated by the period generators.

\subsection{Application}

QFT is a typical component in Quantum algorithms: QFT can be done efficiently with $\mathcal{0}(\catorder^2)$ controlled single-qubit gates.
If the Fourier transform is a basis vector, a single measurement is sufficient to determine it.